\chapter{From design to architecture: Design Principles}

Design principles means follow a clean architecture and a clean architecture means that the 
architecture can easily evolve, to be changed.



\section{Liskov Substitution Principle - LSP}

If we substrute o2 vitw o1 and the behavior of p
does not change means S is derived class of T.

S is a subtype od T.

\section{Interface Seggregation Principle - ISP}

you shoudnt building interfaces or classes with too many operation
that are not essential, not used.

Similart to SRP we have different actor, that represent diff class
U1 use only operation1 and U2 use only operation2 and 
U3 use only operation3

if u wnat to change operation3 you need to redeploy it, 
it is dangerous.
Is diff from the SRP we need to modularize better, how?

Inserting interfaces.The classes use the interfaces instead the
class and the interface implement the class.

if we change operation1 we impact also user2 and user3, if we
change the interface we impact only the corrisponded user.


---

No class with too much code that is not very used in certains 
cases.

\section{Dependecy Inversion Principle - DIP}
You should alwaies have dependecies that lover moduls depends on higher moduls, not the opposite.
Any concrete class only abstract class or interface, not on concrete class.

We can make an interface:

App -> service -> concrete class
App -> iService <- concrete class

We can packeage the app and the interface togheder.

We can do something also the factory from the app -> iFactory <- concrete factory -> concrete class -> iService <- app



  



